\documentclass[aspectratio=169]{beamer}

\usepackage[utf8]{inputenc}
\usepackage[T1]{fontenc}
\usepackage[english]{babel}
\usepackage{lmodern}
\usepackage{hyperref}
\usepackage{booktabs}

\usetheme{Madrid}
\usecolortheme{default}
\setbeamertemplate{navigation symbols}{}

\title{Air Quality Assessment in Torino}
\subtitle{Data Management Project Presentation}
\author{[Student 1] \and [Student 2] \and [Student 3]}
\institute{Universita degli Studi di Milano-Bicocca (UNIMIB)}
\date{Academic Year 2025-2026\\[insert presentation date]}

\begin{document}

\begin{frame}
  \titlepage
\end{frame}

\begin{frame}{Agenda}
\begin{itemize}
  \item Motivation and objectives
  \item Research questions
  \item Data sources and acquisition
  \item Storage and database design
  \item Data quality and integration
  \item Analysis and dashboard
  \item Reproducibility
  \item Limitations and future work
  \item Conclusion
\end{itemize}
\end{frame}

\begin{frame}{Motivation}
\begin{itemize}
  \item Air pollution is a relevant urban and public-health issue.
  \item Torino is a useful case because it includes both Traffic and Background monitoring stations.
  \item The goal is not only to visualize results, but to build a reproducible data-management pipeline.
  \item The project studies PM10, NO2, exceedance days, and integration quality.
\end{itemize}
\end{frame}

\begin{frame}{Objectives}
\begin{itemize}
  \item Acquire heterogeneous datasets automatically.
  \item Store them in a DBMS suitable for analytical querying.
  \item Transform and validate them through a reproducible workflow.
  \item Integrate weather and spatial enrichment into the analytical warehouse.
  \item Answer four research questions through SQL and Streamlit.
\end{itemize}
\end{frame}

\begin{frame}{Research Questions}
\begin{itemize}
  \item RQ1: How do PM10 and NO2 annual concentrations differ between Traffic and Background stations in 2022--2024?
  \item RQ2: How do annual PM10 and NO2 indicators evolve over time?
  \item RQ3: How do PM10 exceedance days above 50 ug/m\textsuperscript{3} vary by year and station type?
  \item RQ4: Can weather and reverse-geocoded station information be integrated with measurable quality?
\end{itemize}
\end{frame}

\begin{frame}{Data Sources}
\begin{columns}[T]
  \column{0.33\textwidth}
  \textbf{EEA DiscoData}
  \begin{itemize}
    \item annual PM10 and NO2 indicators
    \item 669 loaded records
  \end{itemize}

  \column{0.33\textwidth}
  \textbf{Open-Meteo}
  \begin{itemize}
    \item hourly weather observations
    \item 26,328 loaded records
  \end{itemize}

  \column{0.33\textwidth}
  \textbf{Nominatim / OSM}
  \begin{itemize}
    \item reverse geocoding
    \item spatial enrichment and error measurement
  \end{itemize}
\end{columns}
\end{frame}

\begin{frame}{Acquisition Strategy}
\begin{itemize}
  \item Three automated Python dlt pipelines.
  \item API-based ingestion for all sources.
  \item Raw tables loaded into BigQuery.
  \item Repeatable execution through environment variables and cloud credentials.
\end{itemize}

\vspace{0.4cm}
Suggested visual: APIs $\rightarrow$ dlt $\rightarrow$ BigQuery raw.
\end{frame}

\begin{frame}{Storage and Architecture}
\begin{itemize}
  \item Storage layer: Google BigQuery
  \item Transformation layer: dbt
  \item Presentation layer: Streamlit
  \item Query layer: SQL files in the repository
\end{itemize}

\vspace{0.3cm}
Main analytical tables:
\begin{itemize}
  \item fact\_air\_quality\_annual
  \item mart\_station\_type\_year\_summary
  \item mart\_air\_quality\_weather\_yearly
  \item dim\_station\_enriched
  \item dq\_reverse\_geocoding\_quality
\end{itemize}
\end{frame}

\begin{frame}{Data Quality}
\begin{itemize}
  \item Completeness: key fields tested with \texttt{not\_null}
  \item Uniqueness: station keys tested with \texttt{unique}
  \item Validity: accepted station types and enrichment flags
  \item Automated evidence: 30 dbt tests passed
\end{itemize}
\end{frame}

\begin{frame}{Data Integration and Enrichment}
\begin{itemize}
  \item Weather enrichment joins yearly weather indicators to annual air-quality facts.
  \item Spatial enrichment reverse-geocodes monitoring stations.
  \item Integration is fully automated in the warehouse pipeline.
  \item Error is measured explicitly, not only described.
\end{itemize}

\vspace{0.3cm}
Key metrics:
\begin{itemize}
  \item weather enrichment success: 100\%
  \item geocoding coverage: 100\%
  \item average geocoding distance: 27.09 m
  \item maximum geocoding distance: 55.34 m
\end{itemize}
\end{frame}

\begin{frame}{RQ1 and RQ2 Results}
\begin{itemize}
  \item Traffic stations consistently show worse pollutant values than Background stations.
  \item In 2024, NO2 is 38.59 ug/m\textsuperscript{3} for Traffic vs 27.77 ug/m\textsuperscript{3} for Background.
  \item PM10 decreases between 2022 and 2024 for both station types.
  \item NO2 also declines, but remains close to the EU limit at Traffic stations.
\end{itemize}
\end{frame}

\begin{frame}{RQ3 Results}
\begin{itemize}
  \item PM10 exceedance days are higher at Traffic stations in every observed year.
  \item 2022: 71.5 vs 47.5
  \item 2023: 58.0 vs 28.0
  \item 2024: 44.5 vs 39.0
  \item The EU maximum of 35 days is exceeded systematically at Traffic stations.
\end{itemize}
\end{frame}

\begin{frame}{RQ4 Results}
\begin{itemize}
  \item Weather and spatial enrichment were integrated successfully.
  \item Reverse geocoding covered all stations.
  \item Quality remained controlled, with low measured spatial error.
  \item The project satisfies the FAQ requirement on automated enrichment and error measurement.
\end{itemize}
\end{frame}

\begin{frame}{Dashboard}
\begin{itemize}
  \item Streamlit dashboard organized by executive summary and research-question pages.
  \item KPI cards, evidence tables, and presentation-ready charts.
  \item Built directly on BigQuery marts.
  \item Useful both for analysis and for oral presentation support.
\end{itemize}
\end{frame}

\begin{frame}{Reproducibility}
\begin{itemize}
  \item Source code versioned in the repository.
  \item Pipelines implemented in Python.
  \item Warehouse modeled in BigQuery.
  \item Transformations and tests implemented in dbt.
  \item Analytical queries stored as SQL files.
  \item Dashboard launchable through Streamlit.
\end{itemize}
\end{frame}

\begin{frame}{Limitations}
\begin{itemize}
  \item Annual air-quality indicators limit temporal granularity.
  \item The spatial scope is limited to the selected Torino stations.
  \item The analysis is descriptive, not causal.
  \item More external contextual datasets could strengthen interpretation.
\end{itemize}
\end{frame}

\begin{frame}{Future Work}
\begin{itemize}
  \item Add higher-frequency pollution data.
  \item Extend beyond Torino.
  \item Integrate additional spatial datasets such as traffic or land use.
  \item Track data-quality indicators historically.
  \item Add more advanced analytical comparisons and policy context.
\end{itemize}
\end{frame}

\begin{frame}{Conclusion}
\begin{itemize}
  \item End-to-end project: acquisition $\rightarrow$ storage $\rightarrow$ profiling $\rightarrow$ integration $\rightarrow$ analysis $\rightarrow$ quality improvement.
  \item At least two data sources are used, with API-based acquisition.
  \item Integrated data are stored in a DBMS and queried through SQL.
  \item Integration quality is measured explicitly.
  \item The final output includes both analytical results and a reproducible architecture.
\end{itemize}
\end{frame}

\begin{frame}{Questions}
\centering
\Huge Thank you!\\[0.5cm]
\Large Questions?
\end{frame}

\end{document}